\scanpage{23}
\mspage{11}
Écrivons comme tantôt $K=\varinjlim A_i$, donc \deleted{$\unsure{K_0|K}$} $\eta=\operatorname{Spec}(K)=\varprojlim U_i$, on a (en termes d'un pt géom. quelconque $\bar\eta$ de $\operatorname{Spec}K$ i.e. en termes d'un $\bar K$)
\[
\pi_K=\varprojlim_i\pi_1(\bar U_i\deleted{\unsure{\bar\eta|\bar K}},\unsure{\bar\eta|\bar K}),\quad\text{où }\bar U_i=U_i\otimes_k\overline{\mathbb Q}
\]
et il suffit de voir que $\forall i,\ \overset{f_i}{\overline{f|A_i}}=\overset{g_i}{\overline{g|A_i}}$ : moyennant le fait que
\[
\pi_1(f_i^*)=\pi_1(g_i^*):\pi_{K'}\longrightarrow\pi_1(U_i)\deleted{\unsure{K|\bar\eta}}
\]
(comme hom de groupes extérieurs. Or pour $i$ fixé, on a $K'=\varinjlim A'_j$, où les $A'_j$ contiennent $f_i(A_i)$ et $g_i(A_i)$, donc
\[
\pi_{K'}=\varprojlim_j\pi_1(\bar V_j,\bar\eta'),\quad\text{avec }\bar V_j=\operatorname{Spec}(A_j)\otimes_k\overline{\mathbb Q}.
\]
Notons (prenant les $V_j$ \deleted{\unsure{normaux}}\inserted{réguliers}) que les hom de transition du syst. projectif des $\pi_1(\bar V_j,\bar\eta')$ sont surjectifs --- donc $\pi_{K'}\longrightarrow\pi_1(\bar V_j,\bar\eta')$ est surjectif, ce qui implique que l'égalité de $f^*$ et $g^*:\pi_{K'}\rightrightarrows\pi_1\deleted{\unsure{K|k}}(\bar U_i)$ (comme hom extérieurs) implique celle des $\pi_1(\bar V_j)\longrightarrow\pi_1(\bar U_i)$.
Donc l'égalité $f_i=g_i$ (d'où $f=g$) sera consé-quence du résultat plus général,
\margin{22}
